\documentclass[11pt,a4paper]{article}
\usepackage[spanish,es-noshorthands]{babel}
\usepackage[T1]{fontenc}
\usepackage{lmodern}
\usepackage[margin=2.3cm]{geometry}
\usepackage{xcolor}
\usepackage{textcomp}
\usepackage[hidelinks]{hyperref}
\hypersetup{pdfauthor={Ramón Alejandro Manzaba Castro}, pdftitle={Consentimiento informado de la encuesta (versión 2.0)}}
\setlength{\parindent}{0pt}
\setlength{\parskip}{5pt}
\begin{document}
{\small\textit{Texto de la primera pantalla del formulario en línea. Versión 2.0, 25 de noviembre de 2026. Trabajo de Integración Curricular: Metodología para el diagnóstico y gestión de riesgos en el tratamiento de datos personales académicos conforme a la LOPDP en universidades ecuatorianas. Las notas grises indican el numeral del art.~12 de la LOPDP que cubre cada párrafo y no se muestran al participante.}}
\medskip
\begin{center}{\large\bfseries Consentimiento informado: encuesta a estudiantes universitarios del Ecuador (sin datos identificativos directos)}\end{center}


\medskip\noindent\rule{\textwidth}{0.4pt}\medskip

\textbf{Encuesta: Datos personales de estudiantes universitarios del Ecuador}

Le invitamos a participar en una encuesta para un Trabajo de Integración Curricular de la carrera de Ingeniería en Tecnologías de la Información de la Universidad Tecnológica Indoamérica. La encuesta está dirigida a estudiantes de grado de \textbf{todas} las universidades y escuelas politécnicas del Ecuador, públicas y privadas.

\textbf{¿Quién es el responsable de sus datos?} Ramón Alejandro Manzaba Castro, estudiante investigador, bajo la tutoría de María Cristina Jiménez Hernández, MSc. Contacto: ramonalejandromanzaba@gmail.com. Por tratarse de una investigación individual, no existe un delegado de protección de datos. {\footnotesize\color{gray}(LOPDP, art.~12, nums. 8 y 9)}

\textbf{¿Para qué es la encuesta?} Para conocer qué datos personales piden las universidades y escuelas politécnicas a sus estudiantes, si les informan para qué los usan, cómo obtienen su consentimiento cuando hace falta y qué saben los estudiantes sobre sus derechos según la Ley Orgánica de Protección de Datos Personales (LOPDP). Las respuestas no se usarán para ninguna otra finalidad. {\footnotesize\color{gray}(LOPDP, art.~12, nums. 1 y 7)}

\textbf{¿Con qué base legal?} Con su consentimiento, libre, específico, informado e inequívoco (LOPDP, arts. 7, num. 1, y 8), que usted da al marcar "Sí, acepto". {\footnotesize\color{gray}(LOPDP, art.~12, num. 2)}

\textbf{¿Qué se hará con sus respuestas?} Se recogen en un formulario en línea, se guardan en una base de datos de la investigación y se analizan con estadística; los resultados se publican \textbf{solo de forma agregada} (porcentajes del conjunto, nunca respuestas individuales ni resultados de una universidad en particular). No se toman decisiones automatizadas sobre usted ni se elaboran perfiles. {\footnotesize\color{gray}(LOPDP, art.~12, nums. 3, 5 y 17)}

\textbf{¿Quién más tiene acceso? (encargado y transferencia internacional)} Las respuestas se almacenan en el servicio de formularios en línea que se usa para la encuesta (Google Forms o Microsoft Forms; la versión publicada del formulario indica cuál). Ese proveedor actúa como \textbf{encargado del tratamiento} y sus servidores pueden estar \textbf{fuera del Ecuador}, lo que constituye una transferencia internacional. No se comunican las respuestas a ninguna otra persona o institución. {\footnotesize\color{gray}(LOPDP, art.~12, num. 10)}

\textbf{¿Quién puede participar?} Personas de \textbf{18 años o más} que estudien actualmente una carrera de grado en una universidad o escuela politécnica del Ecuador.

\textbf{¿Qué datos se recogen?} \textbf{No le pediremos} su nombre, número de cédula, correo electrónico, número de matrícula ni teléfono; el formulario no registra su correo, no exige iniciar sesión y no permite subir archivos. Le preguntaremos en qué universidad estudia (puede elegir "Prefiero no responder") y, de forma opcional, la región: esos datos se usan solo para el análisis interno y \textbf{nunca se publicará un resultado por universidad}. Las preguntas generales usan rangos (por ejemplo, de edad). No se le preguntará por su salud ni por una posible discapacidad; solo si su universidad le \textbf{pidió} ese tipo de datos. Como algunas respuestas combinadas podrían, en teoría, acercarse a una persona, la base se trata como \textbf{seudonimizada} y no como anónima: en el análisis la universidad se reemplaza por un código y no se cruza con el semestre. Los datos se recogen directamente de usted. {\footnotesize\color{gray}(LOPDP, art.~12, num. 6)}

\textbf{¿Es obligatoria?} No. Su participación es \textbf{voluntaria} y no recibirá compensación. Negarse o dejar preguntas sin responder \textbf{no tiene ninguna consecuencia} para usted; si una respuesta es errónea, solo afecta la calidad del análisis estadístico. Puede abandonar la encuesta en cualquier momento cerrando la página; lo que no envíe no se registrará. {\footnotesize\color{gray}(LOPDP, art.~12, nums. 11 y 12)}

\textbf{¿Puedo retirar mi consentimiento y ejercer mis derechos?} Sí. Puede revocar su consentimiento y ejercer los derechos de acceso, rectificación y actualización, eliminación, oposición, portabilidad, suspensión del tratamiento y a no ser objeto de decisiones automatizadas escribiendo a ramonalejandromanzaba@gmail.com. Como no se recogen datos que le identifiquen directamente, \textbf{es posible que no podamos localizar sus respuestas una vez enviadas}; en ese caso se le explicará por escrito. {\footnotesize\color{gray}(LOPDP, art.~12, nums. 13, 14 y 15)}

\textbf{¿Dónde puedo reclamar?} Ante el investigador, en el correo indicado, y ante la \textbf{Superintendencia de Protección de Datos Personales} (SPDP), autoridad de protección de datos del Ecuador. {\footnotesize\color{gray}(LOPDP, art.~12, num. 16)}

\textbf{¿Cuánto tiempo se conservan?} Al cerrar la encuesta, las respuestas se exportan a un archivo cifrado al que solo accede el investigador y se \textbf{eliminan del servicio de formularios}. La base seudonimizada se conserva hasta la aprobación del Trabajo de Integración Curricular y luego se elimina. Ninguna respuesta individual se publica ni se sube a repositorios. {\footnotesize\color{gray}(LOPDP, art.~12, num. 4)}

\textbf{¿Cuánto dura?} Entre 13 y 15 minutos.

\textbf{¿Hay riesgos o beneficios?} El riesgo es mínimo, porque no se recogen datos identificativos directos ni datos sensibles suyos. El beneficio esperado es contribuir a que las universidades protejan mejor los datos de sus estudiantes.

\textbf{Verificación ética:} en trámite; se completará con la respuesta a la solicitud de verificación ética presentada a la Universidad Tecnológica Indoamérica. La encuesta no se publica sin ella.

\textbf{Recomendación:} en la pregunta abierta final (solo en la prueba piloto), por favor no escriba nombres, números de cédula, correos ni otros datos que le identifiquen a usted o a otras personas.

\medskip\noindent\rule{\textwidth}{0.4pt}\medskip

\textbf{[P00] He leído la información anterior y acepto participar voluntariamente.}
\begin{itemize}
  \item Sí, acepto
  \item No acepto \textit{(la encuesta termina)}
\end{itemize}
\end{document}
